\section{Discussion\label{sec:discussion}}
The research questions examine where architectural signal resides and how a mapping model can use it. RQ1 compares file representations before message passing, RQ2 examines file, member, and folder context and compares the folder gains with a directory-based rule, and RQ3 examines how predictor choice and aggregation direction affect the contribution of dependency neighbors. This section discusses these findings together.

\subsubsection*{Dependency Information in File Representations}
Dependency information can support mapping before message passing, through the graph position encoded by Node2Vec or through dependency-guided adaptation of a code encoder. The advantage of these representations over Word2Vec, and of fine-tuned UniXcoder over its zero-shot counterpart, suggests that dependency information is useful beyond what the evaluated lexical and pretrained code representations provide. The representations also differ in training, vector size, and source content, so these comparisons are between representation pipelines and not between isolated information sources.

The usefulness of this signal varies with the relation between dependency structure and the reference architecture. Node2Vec’s micro-F1 is lower on the three least homophilic systems, SH-3D, HDF, and T.Mates (Table~\ref{tab:datasets}), with a median of $58$ compared with $72$ on the remaining seven. Node2Vec learns dependency proximity from files that co-occur on random walks, without using the GT mapping. When dependencies frequently cross module boundaries, files from different modules may co-occur on walks and receive similar vectors, making the modules harder to distinguish.

The advantage of fine-tuned UniXcoder over Node2Vec on SH-3D and T.Mates suggests that code content may provide complementary evidence in these mixed neighborhoods. HDF does not follow this pattern, so low homophily alone does not explain their relative performance; HDF is also the smallest system, and every representation has comparatively low micro-F1 on it. Homophily is used here to interpret the observed pattern, not to select a representation, since computing it requires the reference mapping.

Encoding dependencies in the initial representation and aggregating dependency neighbors during mapping are two uses of the same information source, and the second does not necessarily replace the first. With undirected aggregation, the median gains are generally larger for Word2Vec and zero-shot UniXcoder than for Node2Vec and fine-tuned UniXcoder, yet fine-tuned UniXcoder still has the highest median macro-F1 for every predictor. Conversely, useful dependency information in a representation does not guarantee a benefit from further aggregation: RQ3 shows that this benefit depends on both aggregation direction and predictor choice.

\subsubsection*{Folder Organization and Architectural Exceptions}
The gains from adding folder nodes in RQ2 show that directory organization helps reproduce the reference mapping. They do not show that this information must reach a file through message passing over folder nodes, since supplying the same information directly in the file vector matches or exceeds the accuracy of folder nodes on most systems. Extending a directory assignment is also different from identifying files that depart from it. Comparing FM, FFM, and FD on the directory baseline's errors helps assess this distinction (Section~\ref{sec:directory_errors}), but does not isolate architectural exceptions from errors caused by missing local supervision.

Lucene and HDF have pure directories under the reference mapping, so their baseline errors do not reflect conflicts between directory and module membership. They may instead arise from directories without labeled examples, although the table does not establish this. FFM recovers most baseline errors on Lucene and half on HDF, but only about a fifth on A.UML, whose directories are also pure. Directory purity therefore does not ensure that folder nodes recover baseline errors.

Bash and HDC have mixed directories. On baseline errors, FFM is similar to FM on Bash and less accurate on HDC. JabRef further illustrates why overall gains are insufficient: adding folder context raises overall accuracy while lowering accuracy on baseline errors.

The directory rule serves as a reference point for interpreting the learned models. Where directory organization closely follows the architecture, it already reproduces most of the mapping. On most evaluated systems, the added modeling complexity therefore brings little accuracy over the rule. Evaluating the additional value of a learned mapper requires examining both files in directories without labeled examples and files whose module differs from their directory's majority module.

\subsubsection*{Dependency Neighbors and Aggregation Direction}

The benefit of dependency neighbors depends on how the predictor aggregates them. With undirected aggregation, HGT gains least, while GAT achieves the highest median macro-F1 with every representation, although differences between predictors are small relative to variation across systems. GAT and HGT differ in their attention formulations and parameterizations, so the comparison does not isolate why GAT performs better. These findings concern the evaluated file-level predictors, not the usefulness of dependency types in general. RQ2 uses HGT for the folder- and member-aware configurations; whether the folder gains persist with other predictors remains untested.

The preferred aggregation direction appears to depend more on the system than on the predictor or representation, although the pattern is not uniform across metrics. Undirected aggregation generally gives larger gains, but several systems consistently favor directed aggregation. Agreement across predictors and representations describes the stability of these system-level preferences rather than providing independent evidence, since the comparisons share systems and splits.

Undirected aggregation adds reverse edges, so the comparison does not isolate propagation direction from neighborhood size. A reversed-edges-only setting would help distinguish these effects and remains for future work.

With Node2Vec, directed aggregation lowers median macro-F1 for all three predictors. Since its vectors were learned on an undirected graph, this may reflect a mismatch between the neighborhoods encoded during initialization and those used during prediction, but this explanation remains untested.

\subsubsection*{Connectivity and Neighbor Information}
Connectivity describes how much neighboring context is available, measured here by average node degree, while homophily describes whether linked files share a module. Few neighbors limit the context available to shallow message passing, whereas Node2Vec can encode more distant context through random walks. Limited connectivity may therefore constrain the benefit of additional aggregation even when the initial structural representation remains useful.

With Node2Vec and undirected aggregation, no predictor changes significantly on Bash and HDC, the two lowest-degree systems. The depth analysis is also consistent with this interpretation: higher average degree is associated with larger gains from a second message-passing layer, particularly in FFM.

Connectivity may also help explain the effect of aggregation direction. T.Mates and SH-3D combine low homophily with high average degree and favor directed aggregation in RQ3. Bash and HDC favor undirected aggregation, as does HDF, which has low homophily but much lower average degree than T.Mates and SH-3D. Reverse edges may supply missing context in sparse neighborhoods but introduce additional cross-module context in dense, low-homophily neighborhoods. This mechanism remains untested: the ten-system comparison cannot separate connectivity from homophily, system size, or label scarcity.

\subsubsection*{Implications for Evaluation}
Evaluation should establish what the input representation already provides, then isolate access to neighborhood information within a fixed predictor. Comparing a GNN model with an MLP changes both the predictor and the use of dependency-neighborhood aggregation. Enabling and disabling neighborhood aggregation within the same predictor family, while keeping the file representations fixed, provides a more targeted comparison. For the neighbor MLP, this comparison includes an increase in input width and parameter count. The same applies to folder information: a comparison of graph variants should include a model that receives that information directly as file features.

Micro-F1 and macro-F1 should be read together, since correctly mapping more files does not necessarily improve performance across modules. For tools intended to identify architectural exceptions, aggregate scores should also be accompanied by performance on files whose module differs from the majority module of their directory.

The intended use also matters. If mapping supports code reorganization, the existing directory structure may be something to reconsider rather than reproduce. Strong performance based on that structure does not establish which information would remain useful when the organization changes. The present study does not test this setting, but motivates evaluating information sources against the intended task rather than treating any representation or predictor as generally preferable.
